% Local report derived from the repository Paper example.
% The upstream ElegantPaper class and local adapter remain in the repository.
\documentclass[lang=en,bibend=bibtex,a4paper,11pt]{elegantpaper}
\usepackage{xurl}
\usepackage{researcher-paper}
\hypersetup{pdftitle={Method A and wider use: an invented comparison}}
\title{Method A and wider use:\\an invented comparison}
\author{}
\date{Scientific version v1 \quad Draft 1}

\begin{document}
\maketitle

\section{Summary}
\label{sec:summary}
Do these inputs justify choosing method A for wider use?
No. The supplied values are invented, and no experiment ran.
They describe A with 8 successes out of 10 and B with 7 out of 10 under the same listed conditions.
That difference describes the example. It does not establish that A performs better in real use.

The input contains one ten-item comparison, with no repeat runs, uncertainty estimate, or external validation.
Human review remains pending, and no dependent scientific action is approved.
Section~\ref{sec:review} asks the researcher to choose the next disposition of this draft.
The source is \href{run:evidence.json}{evidence.json}, version \texttt{fixture-v1}, scientific version \texttt{v1}.
The linked \href{run:evidence-index.json}{evidence index} identifies the exact input and draft with file digests.

\section{Comparison and its reasoning}
\label{sec:comparison}
The question is whether the supplied comparison supports choosing A for wider use.
The input lists two methods under dataset D1, the test split, and seed 7.
A test split is a designated subset for evaluation.
A seed is a number used to initialize randomness.
These are labels in an invented record, not evidence that an evaluation ran.

\begin{table}[htbp]
\centering
\begin{tabular}{lllll}
\toprule
Method & Successes / total & Dataset & Split & Seed \\
\midrule
A & 8 / 10 & D1 & test & 7 \\
B & 7 / 10 & D1 & test & 7 \\
\bottomrule
\end{tabular}
\caption{Invented values under the same listed conditions. A has one more listed success. No experiment ran.}
\label{tab:comparison}
\end{table}

The table reproduces rows E1 and E2 in \href{run:evidence.json}{evidence.json}, at \texttt{/rows/0} and \texttt{/rows/1}.
The totals and condition labels match, so the counts can be compared within this example.
A has one more listed success than B out of the same total of 10.
This is a descriptive difference in the supplied numbers, not a measured advantage.

The shared labels do not establish a fair real-world comparison because the record describes no executed test.
The input also gives no item-level outcomes or definition of success.
The counts therefore supply no supported explanation for the difference.
They do not establish that either method is better for wider use.
This judgment concerns the limits of the supplied input, rather than the actual merits of either method.

\section{Limitations}
\label{sec:limitations}
The source explicitly describes invented inputs for a document-delivery smoke case, a small exercise that tests document production.
It states that no experiment ran.
Only one ten-item comparison appears in the source.
No repeat runs, uncertainty estimate, or external validation support it.
The source locations are \texttt{/provenance} and \texttt{/limitations} in \href{run:evidence.json}{evidence.json}.

An uncertainty estimate describes how imprecise an estimate is.
External validation tests a result beyond its original setting.
Neither is supplied here, and this report adds neither.
The inspected material consists of the complete supplied JSON input, not a paper, implementation, or experiment log.
Real observations under a defined comparison are missing, so the input cannot resolve the choice for wider use.
The report does not prescribe or carry out further scientific work.

\section{Review question}
\label{sec:review}
Does the researcher accept this draft's statement that the invented comparison does not justify choosing A for wider use?
Review the reasoning in Section~\ref{sec:comparison} and the limits in Section~\ref{sec:limitations}.
Table~\ref{tab:comparison} points to the exact rows needed to inspect the numerical comparison.
The review concerns scientific version \texttt{v1}, Draft 1 of \texttt{report.tex}, identified in the \href{run:evidence-index.json}{evidence index}.

The researcher can choose among these dispositions:
\begin{enumerate}
\item Accept this draft as a description of invented inputs and end the writing task. This does not authorize wider use of A.
\item Request changes to the wording or reasoning. The draft remains under review until the researcher decides on its revised version.
\item Define a separate future research task to address the missing real evidence. Its scope and authorization require a separate researcher decision.
\end{enumerate}

The authoritative input field \texttt{/human\_decision} remains \texttt{pending}.
The input field \texttt{/approved\_dependent\_actions} is an empty list.
These fields remain in the unchanged local \href{run:evidence.json}{evidence copy}, and this report records no new approval.
Preparing the draft or building a PDF does not resolve that pending decision.

\end{document}
